% !PRINTED 176
\weberpage{176}
%<<CONT
Theilers $S'$ von $S$ vom Index $j$, aber keine anderen gestattet, so
genügt $\eta$ einer Gleichung $j^{\mathrm{ten}}$ Grades, die als eine Resolvente
des Formenproblems zu betrachten ist. Jede Function, die die
Permutationen von $S'$ gleichfalls gestattet, ist dann eine rationale
Function von $\eta$ und von den Invarianten der Gruppe $S$. Die
Resolvente der $\eta$ ist eine Partial- oder Totalresolvente, je
nachdem die Gruppe $S'$ mit den zu ihr conjugirten Theilern
von $S$ einen gemeinschaftlichen Theiler hat oder relativ prim
ist (Bd. I, \S.\,156).

\begin{center}
\S.\ 44.

\gesperrt{Klein's Erweiterung des algebraischen Grundproblems.}
\end{center}

Die Betrachtungen, die im vorigen Paragraphen durchgeführt
sind, bilden eine directe Verallgemeinerung der \gesperrt{Galois}'schen
Theorie für eine allgemeine Gleichung $n^{\mathrm{ten}}$ Grades; und diese
Theorie ist als Specialfall in der Theorie der linearen Substitutionsgruppen enthalten.

Es sind nämlich nach \S.\ 39 die symmetrischen Functionen
von $n$ unabhängigen Variablen $x_1, x_2, \ldots, x_n$ die Invarianten der
Gruppe $S$, die aus den Permutationen dieser $n$ Variablen besteht,
und die Gleichung $\Phi(t) = 0$ ist also für diesen Fall nach Bd. I,
\S.\ 145 die \gesperrt{Galois}'sche Resolvente der Gleichung $n^{\mathrm{ten}}$ Grades,
deren Wurzeln die Grössen $x_i$ sind. Wir können also die allgemeine Aufgabe der Algebra, eine Gleichung $n^{\mathrm{ten}}$ Grades aufzulösen, als ein Formenproblem einer linearen Substitutionsgruppe
$n^{\mathrm{ter}}$ Dimension auffassen. Nun giebt es aber specielle Gleichungen, die durch Formenprobleme von niedrigerer Dimension
gelöst werden können, so insbesondere die \gesperrt{reinen Gleichungen}
die durch ein Formenproblem der ersten Dimension lösbar sind.

Lineare homogene Substitutionen von einer Dimension sind
nämlich nur von der Form
\weq{1}{x' = \alpha x,}
und wenn diese eine endliche Gruppe vom Grade $\mu$ bilden sollen,
so müssen die Co\"efficienten $\alpha$ Einheitswurzeln vom Grade $\mu$ sein.
Lassen wir umgekehrt $\alpha$ in (1) sämmtliche $\mu^{\mathrm{te}}$ Einheitswurzeln
durchlaufen, so haben wir eine Gruppe vom Grade $\mu$. Diese
Gruppe hat eine absolute Invariante $x^{\mu}$, und wenn diese gegeben
ist, so erhält man $x$ als $\mu^{\mathrm{te}}$ Wurzel daraus.
